\documentclass[12pt,a4paper]{article}

\usepackage[utf8]{inputenc}
\usepackage[T1]{fontenc}
\usepackage{lmodern}
\usepackage{microtype}
% portuges.ldf may be missing in a minimal TeX install; fall back to English
% hyphenation but keep Portuguese captions.
\IfFileExists{portuges.ldf}{%
  \usepackage[brazilian]{babel}%
}{%
  \usepackage[english]{babel}%
  \renewcommand{\contentsname}{Sum\'ario}%
  \renewcommand{\refname}{Refer\^encias}%
  \renewcommand{\figurename}{Figura}%
  \renewcommand{\tablename}{Tabela}%
  \renewcommand{\abstractname}{Resumo}%
}

\usepackage[a4paper,margin=2.5cm]{geometry}
\usepackage{amsmath,amssymb}
\usepackage{graphicx}
\usepackage{booktabs}
\usepackage{longtable}
\usepackage{siunitx}
\usepackage{caption}
\usepackage{subcaption}
\usepackage{float}
\usepackage{xcolor}
\usepackage{listings}
\usepackage{enumitem}
\usepackage[hidelinks]{hyperref}

\graphicspath{{../results/figs/}}

\DeclareSIUnit{\bps}{bps}
\sisetup{group-separator={.},group-minimum-digits=4,output-decimal-marker={,}}

\lstset{
  basicstyle=\ttfamily\small,
  breaklines=true,
  frame=single,
  columns=fullflexible,
  keepspaces=true,
}

\newcommand{\todo}[1]{\textcolor{red}{\textbf{[TODO: #1]}}}
\newcommand{\maybeinput}[1]{%
  \IfFileExists{#1}{\input{\detokenize{#1}}}{\todo{gerar \detokenize{#1}}}}

\title{Compress\~ao PPM-C com Alfabetos de Resolu\c{c}\~ao Vari\'avel\\[2mm]
       \large e Adapta\c{c}\~ao a Fontes N\~ao Estacion\'arias}
\author{Jos\'e Ricardo Rodrigues de Lucena Filho}
\date{Novembro de 2026}

\begin{document}

\maketitle

\begin{abstract}
Este trabalho implementa o algoritmo PPM-C com exclus\~ao para alfabetos de 1, 4 e
8 bits e avalia seu desempenho em fontes n\~ao estacion\'arias do corpus Silesia.
Al\'em do modo sem adapta\c{c}\~ao, o compressor monitora o comprimento m\'edio em
janelas de $j$ s\'imbolos e dispara reset total ou poda de contextos quando a taxa
piora mais de $p\%$. A implementa\c{c}\~ao em C++17 usa codifica\c{c}\~ao
aritm\'etica de 32 bits e tabelas hash por ordem, com verifica\c{c}\~ao bin\'aria
de ida e volta. Em uma grade de calibra\c{c}\~ao de 12 MB (1 MB de cada arquivo),
o A1 teve melhor resultado com reset em $K_{\max}=17$, o A4 com $K_{\max}=8$ e o A8
com $K_{\max}=5$. No corpus completo de 211,9 MB, o reset foi o melhor modo:
3,3546 bits/byte no A1, 2,3628 no A4 e 2,1623 no A8, ganhos de 14,6\%, 10,9\% e
11,9\% sobre o modo sem adapta\c{c}\~ao na mesma configura\c{c}\~ao de cada alfabeto
(A1 em $K_{\max}=17$). As curvas progressivas mostram
estabiliza\c{c}\~ao em poucos megabytes no Dickens e quedas claras nas 11
transi\c{c}\~oes do Silesia, com recupera\c{c}\~ao bem mais curta quando reset ou
poda est\~ao ativos.
O melhor resultado ainda fica cerca de 19\% acima do PPMd de ordem 16 do 7-Zip,
o que reflete a maior complexidade do modelo comercial.
\end{abstract}

\tableofcontents

%======================================================================
\section{Introdu\c{c}\~ao}

\subsection{Contextualiza\c{c}\~ao}
A compress\~ao sem perda explora a redund\^ancia estat\'istica de uma fonte para
representar a mesma informa\c{c}\~ao com menos bits. Modelos baseados em contextos,
como o PPM~\cite{cleary1984} (\emph{Prediction by Partial Matching}), estimam a
probabilidade do pr\'oximo s\'imbolo a partir dos $k$ s\'imbolos anteriores; quanto
maior a ordem $k$, mais espec\'ifica \'e a previs\~ao, mas maior \'e o custo de
armazenamento e o tempo de processamento. O PPM-C~\cite{moffat1990} acrescenta a
esse modelo um mecanismo de
\emph{exclus\~ao}, que evita que um s\'imbolo j\'a considerado em uma ordem superior
seja novamente codificado nas ordens inferiores.

Fontes reais raramente s\~ao estacion\'arias. No corpus Silesia, a concatena\c{c}\~ao
dos doze arquivos cria onze transi\c{c}\~oes bruscas entre tipos de conte\'udo
(texto, c\'odigo, imagens, bases de dados e dados bin\'arios). Um modelo adaptativo
que aprendeu a distribui\c{c}\~ao do arquivo anterior passa a codificar o novo
arquivo com probabilidades inadequadas, e o comprimento m\'edio piora antes de o
modelo se readaptar. Este trabalho estuda esse efeito e avalia duas estrat\'egias
para acelerar a recupera\c{c}\~ao: o rein\'icio total do modelo (\emph{reset}) e a
poda de contextos pouco usados.

\subsection{Fundamenta\c{c}\~ao te\'orica}
O modelo PPM~\cite{bell1990} mant\'em, para cada ordem $k \in \{0,\ldots,K_{\max}\}$, uma tabela que
associa o contexto formado pelos $k$ s\'imbolos mais recentes \`as frequ\^encias dos
s\'imbolos que o sucederam. Para codificar um s\'imbolo, o algoritmo tenta a ordem
mais alta dispon\'ivel. Se o s\'imbolo aparece no contexto, ele \'e codificado com a
distribui\c{c}\~ao desse contexto; caso contr\'ario, emite-se um \emph{escape} e
repete-se o processo na ordem imediatamente inferior. Se todas as ordens falharem,
o s\'imbolo \'e codificado pela ordem $-1$, uniforme sobre o alfabeto.

Neste trabalho usa-se o m\'etodo C de Moffat~\cite{moffat1990} para a probabilidade de escape. Em um
contexto com $n$ ocorr\^encias e $d$ s\'imbolos distintos, a frequ\^encia de escape
\'e $d$, de modo que
\begin{equation}
  P(\text{escape}) = \frac{d}{n+d}.
\end{equation}
Com exclus\~ao, ao escapar de um contexto todos os s\'imbolos observados nesse
contexto s\~ao removidos das ordens inferiores, inclusive da distribui\c{c}\~ao
uniforme de ordem $-1$. As frequ\^encias s\~ao renormalizadas sobre o que resta; se
todos os s\'imbolos do contexto estiverem exclu\'idos, nenhum c\'odigo \'e emitido e
o escape tem probabilidade unit\'aria.

A codifica\c{c}\~ao aritm\'etica~\cite{witten1987} transforma a sequ\^encia de probabilidades em um
\'unico intervalo $[\text{low},\text{high})$ e emite os bits à medida que o
intervalo se estreita. O codificador usa registradores de 32 bits e trata o
\emph{underflow} pelo mecanismo cl\'assico de bits pendentes. Ap\'os cada s\'imbolo,
todos os contextos de $0$ a $K_{\max}$ s\~ao atualizados (sem \emph{update
exclusion}), e as contagens s\~ao reduzidas por $(c+1)/2$ quando o total mais o
n\'umero de distintos ultrapassa 8192.

\subsection{Objetivos}
\begin{itemize}[leftmargin=*]
  \item Implementar PPM-C com exclus\~ao para alfabetos de 1, 4 e 8 bits.
  \item Comparar tr\^es regimes: sem adapta\c{c}\~ao, com reset e com poda.
  \item Medir o efeito de $K_{\max}$, $j$ e $p$ no comprimento m\'edio e no tempo.
  \item Analisar a velocidade de aprendizado (Dickens) e a recupera\c{c}\~ao ap\'os
        transi\c{c}\~oes (Silesia).
  \item Comparar com compressores comerciais (7-Zip: PPMd e LZMA2).
\end{itemize}

%======================================================================
\section{Materiais e M\'etodos}

\subsection{Implementa\c{c}\~ao}
O compressor foi escrito em C++17, sem depend\^encias externas, e organiza-se em
torno de uma classe \texttt{Model} e de um codificador aritm\'etico de 32 bits. Para
cada ordem $k$ h\'a uma tabela hash que associa a chave do contexto (os \'ultimos
$k$ s\'imbolos empacotados em 64 bits) a um n\'o contendo um vetor de pares
(s\'imbolo, contagem) ordenado por s\'imbolo e o total de ocorr\^encias. As
consultas usam busca bin\'aria; as contagens s\~ao esparsas, de modo que um contexto
guarda apenas os s\'imbolos que realmente o sucederam. Um limite global de
contextos (\texttt{mem\_cap}) evita crescimento ilimitado: ao atingi-lo, novos
contextos deixam de ser criados, mas os existentes continuam sendo atualizados.

O arquivo comprimido tem um cabe\c{c}alho de 32 bytes com a assinatura
\texttt{PPMC}, a vers\~ao, o alfabeto, $K_{\max}$, o modo, $j$, $p$, o tamanho
original e o limite de mem\'oria, seguido do fluxo aritm\'etico. O descompressor l\^e
esse cabe\c{c}alho, reconstr\'oi o mesmo modelo e toma as mesmas decis\~oes de
reset e poda, o que garante a sincronia. Os alfabetos A1, A4 e A8 usam,
respectivamente, cada bit, cada nibble e cada byte como s\'imbolo; o \'ultimo
s\'imbolo \'e completado com zeros e o tamanho original em bytes permite truncar a
sa\'ida. Para comparar os tr\^es alfabetos adota-se a m\'etrica comum de
bits por byte, al\'em de bits por s\'imbolo.

A implementa\c{c}\~ao inclui ainda um autoteste (\texttt{ppm selftest}) que verifica
ida e volta para os tr\^es alfabetos, ordens selecionadas at\'e a m\'axima (0, 1, 2,
4 e o valor m\'aximo), arquivos vazios, dados aleat\'orios e os tr\^es modos. Os
experimentos s\~ao disparados pela ferramenta \texttt{ppm bench}, que paraleliza
combina\c{c}\~oes de par\^ametros e grava uma linha por execu\c{c}\~ao em CSV.

\subsection{Monitoramento de taxa local e adapta\c{c}\~ao}
O monitor acompanha o custo estimado de cada s\'imbolo, $-\log_2(\mathrm{freq}/\mathrm{total})$,
calculado em ponto fixo Q16 por meio de uma tabela de logaritmos. O custo \'e
acumulado em janelas consecutivas de $j$ s\'imbolos. O monitor soma o custo do
s\'imbolo na ordem em que ele foi efetivamente codificado; o custo dos escapes n\~ao
\'e contabilizado separadamente. Isso mant\'em a decis\~ao simples e id\^entica nos
dois lados, mas faz o custo monitorado ser menor que o payload real, sobretudo no
A4/A8. Ao fim de cada janela, o custo atual \'e
comparado com o da janela anterior: se
\begin{equation}
  \text{custo atual} > \text{custo anterior} \cdot \left(1 + \frac{p}{100}\right),
\end{equation}
dispara-se o reset ou a poda. Ap\'os um disparo, a janela seguinte n\~ao \'e
comparada com a que motivou a a\c{c}\~ao, pois o modelo acabou de mudar; a
compara\c{c}\~ao reinicia no ciclo seguinte.

Usa-se o custo estimado, e n\~ao o n\'umero de bits efetivamente emitidos, porque ele
depende apenas de decis\~oes j\'a tomadas pelo modelo. Assim, compressor e
descompressor chegam \`a mesma conclus\~ao no mesmo ponto, sem precisar transmitir
nenhum sinal extra. Para inspe\c{c}\~ao, a op\c{c}\~ao \texttt{-{}-monitor-log}
imprime em \texttt{stderr}, tanto na compress\~ao quanto na descompress\~ao, cada
janela (posi\c{c}\~ao $n$, m\'edia estimada e varia\c{c}\~ao percentual em
rela\c{c}\~ao \`a anterior) e o disparo de reset ou poda. Como as duas
execu\c{c}\~oes usam o mesmo custo, as linhas \texttt{[monitor]} coincidem, o que
torna a sincronia verific\'avel na pr\'atica.

\subsection{Reset e poda}
No modo \emph{reset}, todas as tabelas de contexto e o hist\'orico s\~ao limpos e o
modelo volta a aprender do zero, sem os contextos acumulados. No modo \emph{poda},
os contextos vistos uma \'unica vez s\~ao removidos e as contagens dos demais s\~ao
reduzidas por $(c+1)/2$; esse crit\'erio elimina contextos ef\^emeros e faz o modelo
esquecer estat\'isticas antigas sem perder os contextos recorrentes. Em ambos os
casos a decis\~ao \'e tomada apenas com informa\c{c}\~ao dispon\'ivel nos dois lados,
o que mant\'em a sincronia.

\subsection{Metodologia experimental}
Os experimentos usam o corpus Silesia~\cite{silesia}, com 211.938.580 bytes nos
doze arquivos.
Para a grade da se\c{c}\~ao 3.1 montou-se um fluxo de calibra\c{c}\~ao com 1 MB de
cada arquivo (12 MB, contendo as onze transi\c{c}\~oes). A grade cobre A1 com
$K_{\max}$ de 0 a 32 e A4/A8 com $K_{\max}$ de 0 a 8, $j \in
\{1\,000, 5\,000, 10\,000, 25\,000, 50\,000, 100\,000\}$, $p \in
\{5,10,15,20,25,30,40,50\}\%$ e os tr\^es modos, totalizando 3.201 combina\c{c}\~oes
para A1 e 873 para cada um dos outros alfabetos. Nas tabelas, os modos aparecem
com os nomes usados no c\'odigo: \emph{none} (sem adapta\c{c}\~ao), \emph{reset} e
\emph{prune} (poda).

Para cada execu\c{c}\~ao registram-se o comprimento m\'edio do payload
(bits/s\'imbolo e bits/byte), o tempo de compress\~ao, o tempo de descompress\~ao
(quando a verifica\c{c}\~ao \'e solicitada), o n\'umero de resets e de podas e o
resultado da compara\c{c}\~ao bin\'aria entre original e descomprimido. Cada
compress\~ao usa uma thread; como a grade executa v\'arias em paralelo, os tempos da
grade incluem conten\c{c}\~ao e servem para comparar ordens, n\~ao como tempos
absolutos. As configura\c{c}\~oes escolhidas para cada alfabeto e modo s\~ao depois
validadas no corpus completo, e os resultados s\~ao comparados com o 7-Zip (PPMd
com ordens 6 e 16 e LZMA2). A grade usa um limite de 8 milh\~oes de contextos; a
valida\c{c}\~ao final usa 20 milh\~oes, para n\~ao truncar as ordens mais altas no
corpus completo.

Os experimentos foram executados em uma m\'aquina Linux com 16 n\'ucleos e 31 GB de
mem\'oria. O tempo reportado em cada linha \'e o tempo de parede daquela execu\c{c}\~ao
individual, incluindo a conten\c{c}\~ao com as demais.

%======================================================================
\section{Resultados}

\subsection{An\'alise de ordem e desempenho (3.1)}
\maybeinput{tables/31_A1_kmax.tex}
\maybeinput{tables/31_A4_kmax.tex}
\maybeinput{tables/31_A8_kmax.tex}
\maybeinput{tables/31_summary.tex}
\maybeinput{tables/final.tex}
\maybeinput{tables/dec_times.tex}

\begin{figure}[H]
  \centering
  \begin{subfigure}{.32\linewidth}
    \includegraphics[width=\linewidth]{bpb_vs_kmax_A1.png}
    \caption{A1}
  \end{subfigure}\hfill
  \begin{subfigure}{.32\linewidth}
    \includegraphics[width=\linewidth]{bpb_vs_kmax_A4.png}
    \caption{A4}
  \end{subfigure}\hfill
  \begin{subfigure}{.32\linewidth}
    \includegraphics[width=\linewidth]{bpb_vs_kmax_A8.png}
    \caption{A8}
  \end{subfigure}
  \caption{Bits/byte em fun\c{c}\~ao de $K_{\max}$, usando o melhor $(j,p)$ de cada
  modo no fluxo de calibra\c{c}\~ao.}
  \label{fig:kmax}
\end{figure}

As Tabelas~\ref{tab:kmax-A1}--\ref{tab:kmax-A8} trazem as duas m\'etricas,
bits/byte e bits/s\'imbolo, para cada $K_{\max}$ e modo; a
Tabela~\ref{tab:resumo-31} tamb\'em apresenta ambas. A grade completa (4.947
linhas) est\'a em
\texttt{results/31\_A*.csv}; as Tabelas~\ref{tab:kmax-A1}--\ref{tab:kmax-A8}
resumem, para cada $K_{\max}$, o melhor $(j,p)$ de cada modo.

No alfabeto A1, o modo sem adapta\c{c}\~ao continua melhorando lentamente com a
ordem e atinge o melhor valor da grade em $K_{\max}=30$ (4,270 bits/byte). Os modos
adaptativos t\^em comportamento diferente: tanto o reset quanto a poda t\^em \'optimo
em $K_{\max}=17$ (3,899 e 4,093 bits/byte, respectivamente) e pioram para ordens
maiores. Isso ocorre porque, acima desse ponto, os contextos tornam-se quase
\'unicos e o modelo passa a memorizar o passado em vez de generalizar; al\'em disso,
cada n\'ivel adicional custa tempo e mem\'oria. O tempo m\'edio por combina\c{c}\~ao
na grade paralela do A1 cresce de cerca de 6,5 s em $K_{\max}=0$ para 86 s em
$K_{\max}=17$ e 790 s em $K_{\max}=32$; esses valores incluem a conten\c{c}\~ao entre
as execu\c{c}\~oes simult\^aneas e servem para comparar ordens, n\~ao como tempo
absoluto. A Tabela~\ref{tab:dec-times} repete as melhores configura\c{c}\~oes
isoladamente.

No A4, o melhor resultado sem adapta\c{c}\~ao aparece em $K_{\max}=7$ (3,255
bits/byte), enquanto reset e poda preferem $K_{\max}=8$ (3,018 e 3,080 bits/byte).
No A8, todos os modos t\^em \'optimo em $K_{\max}=5$ (2,992, 2,854 e 2,912
bits/byte). A diferen\c{c}a entre alfabetos \'e coerente com o n\'umero de
s\'imbolos: o A8 precisa de menos ordens para cobrir os padr\~oes de byte, enquanto o
A1, com apenas dois s\'imbolos, suporta ordens muito maiores antes de saturar.

O efeito de $j$ \'e mais forte que o de $p$. Janelas de 1.000 s\'imbolos s\~ao
claramente ruins: no A1 com reset, o melhor resultado com $j=1.000$ \'e 5,353
bits/byte, contra 3,899 no \'otimo ($j=50.000$); no A1 com poda, 6,044 contra
4,093. A partir de 10.000 s\'imbolos o custo j\'a fica pr\'oximo do m\'inimo; no A4
e no A8, 5.000 j\'a \'e suficiente. Janelas de 50.000--100.000 n\~ao trazem ganho
adicional e podem atrasar a rea\c{c}\~ao.
Para $p$, a curva \'e mais plana: no A1 com reset, $p=5$ d\'a 4,186 e $p=30$ d\'a
3,899, enquanto $p=40$--$50$ volta a 3,932--3,921. No A4 e no A8, os melhores
resultados usam $p=40$--$50$ e $j=5.000$--$10.000$, coerente com as transi\c{c}\~oes
abruptas do fluxo de calibra\c{c}\~ao. No A4 com reset, $j=1.000$ d\'a 3,322 contra
3,018 no \'otimo; no A8, 2,957 contra 2,854. A poda \'e mais sens\'ivel a $j$ que o reset
(A1: 6,044 em $j=1.000$ contra 4,093 no \'otimo), porque cada disparo varre as
tabelas; por isso seu melhor $j$ no A1 \'e 50.000, com $p=15$.

A Tabela~\ref{tab:resumo-31} resume as melhores configura\c{c}\~oes da
calibra\c{c}\~ao e compara cada modo com o melhor \emph{none} de cada alfabeto (A1
em $K_{\max}=30$, A4 em 7 e A8 em 5). No corpus completo, para comparar os modos
sem confundir o efeito da ordem, usa-se a mesma configura\c{c}\~ao
($K_{\max},j,p$) nos tr\^es modos de cada alfabeto (A1 em $K_{\max}=17$, A4 em 8 e
A8 em 5). Nessa compara\c{c}\~ao controlada, o reset foi o melhor modo nos tr\^es
casos: 3,3546 bits/byte no A1, 2,3628 no A4 e 2,1623 no A8, contra 3,9260, 2,6522
e 2,4538 sem adapta\c{c}\~ao. Isso corresponde a ganhos de 14,6\%, 10,9\% e 11,9\%
na mesma configura\c{c}\~ao. A poda tamb\'em melhora os tr\^es alfabetos, mas fica atr\'as do
reset (3,4382, 2,5250 e 2,2403 bits/byte). A Tabela~\ref{tab:final} traz os tempos
de compress\~ao e descompress\~ao medidos nessa valida\c{c}\~ao. A grade registra o
tempo de compress\~ao de cada combina\c{c}\~ao; a Tabela~\ref{tab:dec-times} repete
as nove melhores configura\c{c}\~oes em execu\c{c}\~ao isolada, com descompress\~ao.
Medir a descompress\~ao de todas as 4.947 combina\c{c}\~oes dobraria o custo da
grade.

A Tabela~\ref{tab:dec-times} mostra que compress\~ao e descompress\~ao t\^em custos
pr\'oximos, com raz\~ao desc/comp entre 0,97 e 1,15. O caso mais caro \'e o A1 sem
adapta\c{c}\~ao em $K_{\max}=30$ (88,8 s e 102,0 s), enquanto o reset em
$K_{\max}=17$ cai para 16,5 s e 16,0 s; o A8 \'e o mais r\'apido (6,5--9,5 s). Isso
confirma que ordens altas no A1 s\~ao caras tanto para comprimir quanto para
descomprimir; o reset em $K_{\max}=17$, al\'em de comprimir melhor, tamb\'em \'e
muito mais barato que o modo sem adapta\c{c}\~ao em $K_{\max}=30$.

Cabe registrar uma escolha metodol\'ogica: para comparar os modos sem confundir o
efeito da ordem, a valida\c{c}\~ao final usa a mesma configura\c{c}\~ao
($K_{\max}$, $j$ e $p$) nos tr\^es modos de cada alfabeto. No A1, o melhor ponto
isolado do modo sem adapta\c{c}\~ao na grade \'e $K_{\max}=30$ (4,270 bits/byte),
mas o \'optimo global (reset) est\'a em 17. Usar 17 em todos os modos mant\'em a
compara\c{c}\~ao controlada; por isso o ganho de 14,6\% \'e relativo ao \emph{none}
em $K_{\max}=17$, e n\~ao ao melhor \emph{none}. Na calibra\c{c}\~ao, o \emph{none}
em $K_{\max}=30$ \'e 8,8\% melhor que em 17, de modo que o ganho contra o melhor
\emph{none} do A1 \'e menor que 14,6\%. A conclus\~ao qualitativa de que o reset \'e superior se
mant\'em, mas o n\'umero deve ser lido como uma compara\c{c}\~ao na mesma
configura\c{c}\~ao. A poda, por varrer todas as tabelas a cada disparo, foi mais
lenta que o reset nos tr\^es alfabetos, com o pior caso no A4 (1.243 s de
compress\~ao contra 458 s do reset).

Para completar a lacuna de tempos por $K_{\max}$, as
Tabelas~\ref{tab:times-A1}--\ref{tab:times-A8} repetem, para cada $K_{\max}$ e
modo, a melhor $(j,p)$ da grade em execu\c{c}\~ao isolada, registrando compress\~ao e
descompress\~ao. Elas mostram o crescimento do custo com a ordem e confirmam que,
no A1, o reset em $K_{\max}=17$ \'e muito mais barato que as ordens altas do modo
sem adapta\c{c}\~ao. Pequenas diferen\c{c}as em rela\c{c}\~ao \`a Tabela~\ref{tab:dec-times}
s\~ao ru\'ido de medi\c{c}\~ao entre execu\c{c}\~oes independentes.

\maybeinput{tables/times_A1_kmax.tex}
\maybeinput{tables/times_A4_kmax.tex}
\maybeinput{tables/times_A8_kmax.tex}

\subsection{Aprendizado no Collected Works of Charles Dickens (3.2)}
\begin{figure}[H]
  \centering
  \begin{subfigure}{.32\linewidth}
    \includegraphics[width=\linewidth]{trace_dickens_A1.png}
    \caption{A1}
  \end{subfigure}\hfill
  \begin{subfigure}{.32\linewidth}
    \includegraphics[width=\linewidth]{trace_dickens_A4.png}
    \caption{A4}
  \end{subfigure}\hfill
  \begin{subfigure}{.32\linewidth}
    \includegraphics[width=\linewidth]{trace_dickens_A8.png}
    \caption{A8}
  \end{subfigure}
  \caption{Comprimento m\'edio progressivo no arquivo \texttt{dickens}, em
  bits/s\'imbolo (bits efetivamente emitidos). Para comparar com bits/byte,
  multiplicar A1 por 8 e A4 por 2 (A8 j\'a est\'a em bits/byte).}
  \label{fig:dickens}
\end{figure}

As tr\^es curvas come\c{c}am com comprimento m\'edio alto e caem rapidamente nos
primeiros megabytes de dados, \`a medida que o modelo aprende os padr\~oes
locais. Adotando uma toler\^ancia de 5\% em torno do valor final, as curvas
permanecem dentro da faixa a partir de cerca de 0,35 MB (A1), 2,77 MB (A4) e
2,99 MB (A8), ou seja, 3,4\%,
27,2\% e 29,3\% do arquivo. Com toler\^ancia de 1\%, a estabiliza\c{c}\~ao ocorre
por volta de 2,46 MB no A1, 7,11 MB no A4 e 7,51 MB no A8. O A1 converge muito mais
r\'apido; A4 e A8 precisam de v\'arios megabytes porque o alfabeto maior tem mais
s\'imbolos e mais escapes at\'e o modelo cobrir os padr\~oes relevantes.

Em bits/byte, os valores finais s\~ao 3,0213 (A1), 2,0505 (A4) e 1,9075 (A8).
O A1 precisa de mais bits por byte porque cada s\'imbolo \'e um bit e o modelo
paga o custo de codificar uma decis\~ao bin\'aria; o A8 aproveita diretamente os
padr\~oes de byte. As curvas mostram os bits efetivamente emitidos pelo
codificador aritm\'etico. O monitor de reset/poda usa uma estimativa separada,
$-\log_2(\mathrm{freq}/\mathrm{total})$, que n\~ao inclui o custo dos escapes; isso
afeta as decis\~oes de adapta\c{c}\~ao, mas n\~ao o eixo das curvas.

\subsection{Aprendizado e estacionariedade no Silesia (3.3)}
\begin{figure}[H]
  \centering
  \begin{subfigure}{.32\linewidth}
    \includegraphics[width=\linewidth]{trace_silesia_A1_none.png}
    \caption{sem adapta\c{c}\~ao}
  \end{subfigure}\hfill
  \begin{subfigure}{.32\linewidth}
    \includegraphics[width=\linewidth]{trace_silesia_A1_reset.png}
    \caption{reset}
  \end{subfigure}\hfill
  \begin{subfigure}{.32\linewidth}
    \includegraphics[width=\linewidth]{trace_silesia_A1_prune.png}
    \caption{poda}
  \end{subfigure}
  \caption{Comprimento m\'edio progressivo no fluxo concatenado do Silesia,
  alfabeto A1. As linhas tracejadas marcam as 11 transi\c{c}\~oes entre arquivos.}
  \label{fig:silesia-A1}
\end{figure}

\begin{figure}[H]
  \centering
  \begin{subfigure}{.32\linewidth}
    \includegraphics[width=\linewidth]{trace_silesia_A4_none.png}
    \caption{sem adapta\c{c}\~ao}
  \end{subfigure}\hfill
  \begin{subfigure}{.32\linewidth}
    \includegraphics[width=\linewidth]{trace_silesia_A4_reset.png}
    \caption{reset}
  \end{subfigure}\hfill
  \begin{subfigure}{.32\linewidth}
    \includegraphics[width=\linewidth]{trace_silesia_A4_prune.png}
    \caption{poda}
  \end{subfigure}
  \caption{Comprimento m\'edio progressivo no fluxo concatenado do Silesia,
  alfabeto A4, com as 11 transi\c{c}\~oes marcadas.}
  \label{fig:silesia-A4}
\end{figure}

\begin{figure}[H]
  \centering
  \begin{subfigure}{.32\linewidth}
    \includegraphics[width=\linewidth]{trace_silesia_A8_none.png}
    \caption{sem adapta\c{c}\~ao}
  \end{subfigure}\hfill
  \begin{subfigure}{.32\linewidth}
    \includegraphics[width=\linewidth]{trace_silesia_A8_reset.png}
    \caption{reset}
  \end{subfigure}\hfill
  \begin{subfigure}{.32\linewidth}
    \includegraphics[width=\linewidth]{trace_silesia_A8_prune.png}
    \caption{poda}
  \end{subfigure}
  \caption{Comprimento m\'edio progressivo no fluxo concatenado do Silesia,
  alfabeto A8, com as 11 transi\c{c}\~oes marcadas.}
  \label{fig:silesia-A8}
\end{figure}

As curvas mostram quedas claras de desempenho nas transi\c{c}\~oes entre os doze
arquivos; como em 3.2, usam-se os bits efetivamente emitidos. Medindo o custo local
em janelas equivalentes a 2 MB antes e depois de
cada transi\c{c}\~ao, o custo posterior \'e, em m\'edia, 1,65 vezes o
anterior no A1 sem adapta\c{c}\~ao; com reset e poda essa raz\~ao cai para 1,38 e
1,40. No A4, a raz\~ao cai de 1,77 para 1,63 (reset) e 1,60 (poda). No A8 a
diferen\c{c}a entre os modos \'e menor (1,79 contra 1,75 e 1,77), mas o ganho
acumulado no arquivo inteiro ainda \'e de 11,9\% com reset, porque os efeitos se
somam ao longo das doze partes.

O tempo de recupera\c{c}\~ao confirma o benef\'icio. Tomando como refer\^encia o
custo local 10\% acima do n\'ivel anterior \`a transi\c{c}\~ao, a mediana de
recupera\c{c}\~ao no A1 cai de 1,20 milh\~ao de s\'imbolos sem adapta\c{c}\~ao para
0,15 milh\~ao com reset e 0,32 milh\~ao com poda. No A4, cai de 0,27 milh\~ao para
0,10 e 0,08 milh\~ao. No A8 a mediana j\'a \'e pequena nos tr\^es casos
(cerca de 0,09--0,11 milh\~ao), mas a m\'edia melhora de 4,96 milh\~oes para 3,53
milh\~oes (reset) e 3,74 milh\~oes (poda). O reset disparou 973 vezes no A1, 886 no
A4 e 505 no A8; a poda, 980, 933 e 464 vezes. O reset \'e mais agressivo e
restabelece rapidamente um modelo limpo; a poda preserva contextos recorrentes, mas
paga o custo de varrer as tabelas.

No acumulado, o reset foi o melhor modo nos tr\^es alfabetos (3,3546, 2,3628 e
2,1623 bits/byte), seguido pela poda (3,4382, 2,5250 e 2,2403) e pelo modo sem
adapta\c{c}\~ao (3,9260, 2,6522 e 2,4538). No A4 a poda ficou mais distante do
reset (6,9\% acima, contra 2,5\% no A1 e 3,6\% no A8) e foi a mais penalizada em
tempo, por causa das varreduras completas.

\subsection{Integridade}
As nove configura\c{c}\~oes finais (tr\^es alfabetos e tr\^es modos) foram executadas
com a op\c{c}\~ao \texttt{--verify}, que descomprime o payload e compara o resultado
com o original. Todas as linhas de \texttt{final\_A*.csv} terminam em
\texttt{verify=ok}, e as nove execu\c{c}\~oes da Tabela~\ref{tab:dec-times} tamb\'em
foram verificadas com sucesso. Al\'em disso, o autoteste (\texttt{ppm selftest}) cobre os tr\^es
alfabetos, ordens selecionadas at\'e a m\'axima (0, 1, 2, 4 e o valor m\'aximo),
arquivos vazios, dados aleat\'orios e os tr\^es modos. Como demonstra\c{c}\~ao
expl\'icita, o bloco abaixo mostra o caso A8 com poda; o mesmo procedimento foi
repetido para A1 (sem adapta\c{c}\~ao) e A4 (reset) no arquivo \texttt{dickens}:

\begin{lstlisting}
$ ./build/ppm c data/silesia/dickens /tmp/dickens_A8.ppm \
      --bits 8 --kmax 5 --mode prune --j 10000 --p 50
$ ./build/ppm d /tmp/dickens_A8.ppm /tmp/dickens_A8.out
$ cmp data/silesia/dickens /tmp/dickens_A8.out && echo OK
OK
\end{lstlisting}

Os tr\^es testes de ida e volta em \texttt{dickens} terminaram com \texttt{cmp}
retornando sucesso (A1 sem adapta\c{c}\~ao, A4 com reset e A8 com poda), al\'em do
autoteste completo.

\subsection{Compara\c{c}\~ao com softwares comerciais}
\maybeinput{tables/commercial.tex}

O 7-Zip~\cite{sevenzip} foi usado como refer\^encia, com PPMd nas ordens 6 e 16 (256 MB de
mem\'oria) e LZMA2, todos em uma thread. No corpus completo, o melhor resultado
comercial foi o PPMd de ordem 16, com 1,8132 bits/byte, seguido por LZMA2 (1,8377)
e pelo PPMd de ordem 6 (1,8831). O melhor resultado deste trabalho foi 2,1623
bits/byte (A8 com reset), cerca de 19\% acima do PPMd de ordem 16. A diferen\c{c}a
tem origem no modelo: o PPMd do 7-Zip usa ordem 16 sobre bytes, mem\'oria de 256 MB
e uma implementa\c{c}\~ao muito otimizada, enquanto este trabalho usa ordens menores
(A8 com $K_{\max}=5$, A1 com 17) e uma estrutura de hash simples, sem SSE nem
\emph{update exclusion}.

No arquivo \texttt{dickens}, o PPMd de ordem 16 atinge 1,7893 bits/byte e o LZMA2
2,2221; o A8 sem adapta\c{c}\~ao deste trabalho fica em 1,9075 bits/byte, coerente
com o valor final da curva 3.2. Ou seja, em um arquivo pequeno e
homog\^eneo a diferen\c{c}a \'e bem menor do que no corpus completo. Os tempos do
7-Zip n\~ao foram medidos; al\'em disso, eles n\~ao seriam diretamente compar\'aveis,
pois os compressores t\^em implementa\c{c}\~oes, threads e estruturas de dados
diferentes. A compara\c{c}\~ao aqui \'e de comprimento m\'edio.

%======================================================================
\section{Discuss\~ao}
A principal dificuldade de implementa\c{c}\~ao foi garantir que compressor e
descompressor tomassem exatamente as mesmas decis\~oes de reset e poda. A solu\c{c}\~ao
foi basear o monitor no custo estimado $-\log_2(\mathrm{freq}/\mathrm{total})$ do s\'imbolo escolhido,
que depende apenas de decis\~oes j\'a tomadas, e n\~ao no n\'umero de bits
efetivamente emitidos pelo codificador aritm\'etico. O monitor n\~ao soma o custo dos
escapes nem o arredondamento do codificador; por isso o custo estimado \'e menor que
o payload real (cerca de 0,2\% no A1, 3\% no A4 e 7\% no A8 no Dickens, no modo sem
adapta\c{c}\~ao), mas a decis\~ao \'e id\^entica nos dois lados e mant\'em a sincronia.
As curvas de 3.2 e 3.3, por outro lado, usam os bits efetivamente emitidos; a
estimativa acima serve apenas para as decis\~oes de reset/poda.

O custo computacional foi o fator mais limitante. A grade completa do A1 consumiu
cerca de 226 horas de CPU, contra aproximadamente 5 horas para cada um dos outros
dois alfabetos. Ordens altas no A1 criam contextos quase \'unicos, o que n\~ao s\'o
n\~ao melhora a compress\~ao como tamb\'em aumenta muito o tempo e o uso de
mem\'oria; o limite de 8 milh\~oes de contextos usado na grade e o crescimento
aproximadamente linear do n\'umero de tabelas explicam a satura\c{c}\~ao observada. Por isso a grade
de 3.1 foi feita no fluxo de calibra\c{c}\~ao de 12 MB, e as melhores configura\c{c}\~oes
foram validadas no corpus completo. Essa \'e uma limita\c{c}\~ao metodol\'ogica que
deve ser lida com clareza: a grade n\~ao cobre todas as combina\c{c}\~oes no corpus
de 212 MB, o que seria proibitivo, mas a valida\c{c}\~ao final confirma as conclus\~oes
principais.

Outro ponto cr\'itico \'e o custo da poda. Como ela varre todas as tabelas a cada
disparo, seu tempo cresce com o n\'umero de contextos e ficou bem acima do reset no
A4 (1.243 s contra 458 s para comprimir o corpus completo), embora a compress\~ao
resultante seja pr\'oxima. O reset \'e mais simples e mais r\'apido, mas descarta
todo o aprendizado; a poda preserva contextos recorrentes ao pre\c{c}o de varreduras
peri\'odicas. Neste corpus, o reset foi a melhor escolha global nos tr\^es alfabetos.

A compara\c{c}\~ao com o 7-Zip mostra que ainda h\'a espa\c{c}o para melhorar. O PPMd
de ordem 16, com 256 MB de mem\'oria, atinge 1,8132 bits/byte no corpus completo,
enquanto o melhor resultado deste trabalho foi 2,1623 (A8 com reset), uma diferen\c{c}a
de cerca de 19\%. Parte da diferen\c{c}a vem da ordem e do tamanho do modelo, parte
vem de t\'ecnicas que n\~ao foram implementadas, como \emph{update exclusion},
\emph{secondary symbol estimation} (SSE) e uma pol\'itica de mem\'oria mais
eficiente. O compressor em si \'e single-thread (a grade apenas paraleliza
configura\c{c}\~oes independentes) e usa tabelas hash com
\texttt{std::unordered\_map} e vetores ordenados, o que tamb\'em limita a velocidade
em ordens altas.

Entre as limita\c{c}\~oes, destacam-se: a calibra\c{c}\~ao em 12 MB, os tempos de
descompress\~ao medidos para as melhores configura\c{c}\~oes (calibra\c{c}\~ao e corpus
completo), mas n\~ao para todas as linhas da grade, o uso de um custo estimado
que n\~ao contabiliza os escapes, o limite fixo de contextos e a
aus\^encia de um mecanismo espec\'ifico para detectar fronteiras de arquivo. Ainda
assim, os objetivos de implementar PPM-C com exclus\~ao, comparar os tr\^es regimes e
analisar a aprendizagem e a n\~ao estacionariedade foram cumpridos.

%======================================================================
\section{Conclus\~ao}
Este trabalho implementou um compressor e descompressor PPM-C com exclus\~ao para
alfabetos de 1, 4 e 8 bits, com monitoramento de taxa local e estrat\'egias de reset
e poda sincronizadas entre compressor e descompressor. A implementa\c{c}\~ao passou
no autoteste e na verifica\c{c}\~ao bin\'aria das nove configura\c{c}\~oes finais.

Na calibra\c{c}\~ao, o melhor resultado do A1 foi obtido com reset em
$K_{\max}=17$ (3,899 bits/byte), enquanto o A4 e o A8 preferiram $K_{\max}=8$ e
$K_{\max}=5$. No corpus completo, o reset foi o melhor modo nos tr\^es alfabetos:
3,3546 bits/byte no A1, 2,3628 no A4 e 2,1623 no A8, com ganhos de 14,6\%, 10,9\% e
11,9\% sobre o modo sem adapta\c{c}\~ao na mesma configura\c{c}\~ao de cada alfabeto.
No A1, esse ganho \'e uma compara\c{c}\~ao controlada em $K_{\max}=17$, n\~ao contra o
melhor \emph{none} da calibra\c{c}\~ao ($K_{\max}=30$); o ganho contra esse
melhor \emph{none} \'e menor. A poda tamb\'em trouxe ganhos, mas foi mais
lenta e ficou atr\'as do reset em compress\~ao.

As curvas de aprendizagem mostram que o modelo se aproxima do regime est\'avel em
0,35 MB no A1 e em cerca de 2,8--3,0 MB no A4/A8 (toler\^ancia de 5\%) e que as
transi\c{c}\~oes do Silesia produzem degrada\c{c}\~ao imediata, com recupera\c{c}\~ao
muito mais r\'apida quando o reset ou a poda est\~ao ativos. O reset reduziu a
mediana de recupera\c{c}\~ao no A1 de 1,20 milh\~ao para 0,15 milh\~ao de s\'imbolos.
Em compara\c{c}\~ao com o 7-Zip PPMd de ordem 16 (1,8132 bits/byte no corpus
completo), o melhor resultado deste trabalho ficou cerca de 19\% acima, o que \'e
coerente com um modelo mais simples e com menos mem\'oria.

Como trabalhos futuros, podem-se explorar ordens adaptativas por alfabeto,
\emph{update exclusion}, SSE, uma pol\'itica de poda mais barata e uma
implementa\c{c}\~ao paralela, al\'em de repetir a grade completa diretamente no
corpus de 212 MB com mais recursos computacionais.

%======================================================================
\begin{thebibliography}{9}
\bibitem{moffat1990}
A.~Moffat.
\newblock Implementing the PPM data compression scheme.
\newblock \emph{IEEE Transactions on Communications}, 38(11):1917--1921, 1990.

\bibitem{cleary1984}
J.~G. Cleary and I.~H. Witten.
\newblock Data compression using adaptive coding and partial string matching.
\newblock \emph{IEEE Transactions on Communications}, 32(4):396--402, 1984.

\bibitem{witten1987}
I.~H. Witten, R.~M. Neal, and J.~G. Cleary.
\newblock Arithmetic coding for data compression.
\newblock \emph{Communications of the ACM}, 30(6):520--540, 1987.

\bibitem{bell1990}
T.~C. Bell, J.~G. Cleary, and I.~H. Witten.
\newblock \emph{Text Compression}.
\newblock Prentice Hall, 1990.

\bibitem{silesia}
S.~Deorowicz.
\newblock Silesia compression corpus.
\newblock \url{https://sun.aei.polsl.pl/~sdeor/index.php?page=silesia}.

\bibitem{sevenzip}
I.~Pavlov.
\newblock 7-Zip file archiver.
\newblock \url{https://www.7-zip.org/}.
\end{thebibliography}

\end{document}
